\section{AMC 8 1989}
        \begin{enumerate}
        	\item (\textit{AMC 8 1989}) $(1+11+21+31+41)+(9+19+29+39+49)=$
$\text{(A)}\ 150 \qquad \text{(B)}\ 199 \qquad \text{(C)}\ 200 \qquad \text{(D)}\ 249 \qquad \text{(E)}\ 250$


	\item (\textit{AMC 8 1989}) $\frac{2}{10}+\frac{4}{100}+\frac{6}{1000} =$
$\text{(A)}\ .012 \qquad \text{(B)}\ .0246 \qquad \text{(C)}\ .12 \qquad \text{(D)}\ .246 \qquad \text{(E)}\ 246$


	\item (\textit{AMC 8 1989}) Which of the following numbers is the largest?


 $\text{(A)}\ .99 \qquad \text{(B)}\ .9099 \qquad \text{(C)}\ .9 \qquad \text{(D)}\ .909 \qquad \text{(E)}\ .9009$


	\item (\textit{AMC 8 1989}) Estimate to determine which of the following numbers is closest to $\frac{401}{.205}$.


 $\text{(A)}\ .2 \qquad \text{(B)}\ 2 \qquad \text{(C)}\ 20 \qquad \text{(D)}\ 200 \qquad \text{(E)}\ 2000$


	\item (\textit{AMC 8 1989}) $-15+9\times (6\div 3) =$
$\text{(A)}\ -48 \qquad \text{(B)}\ -12 \qquad \text{(C)}\ -3 \qquad \text{(D)}\ 3 \qquad \text{(E)}\ 12$


	\item (\textit{AMC 8 1989}) If the markings on the number line are equally spaced, what is the number $\text{y}$?


 
                    \begin{figure}[H]
                        \centering
                        \includegraphics[width=0.5\textwidth]{Problem/AMC_8_AJHSME_1989_6_Problem_Figure_1.png}
                    \end{figure}
                    
$\text{(A)}\ 3 \qquad \text{(B)}\ 10 \qquad \text{(C)}\ 12 \qquad \text{(D)}\ 15 \qquad \text{(E)}\ 16$


	\item (\textit{AMC 8 1989}) If the value of $20$ quarters and $10$ dimes equals the value of $10$ quarters and $n$ dimes, then $n=$
$\text{(A)}\ 10 \qquad \text{(B)}\ 20 \qquad \text{(C)}\ 30 \qquad \text{(D)}\ 35 \qquad \text{(E)}\ 45$


	\item (\textit{AMC 8 1989}) $(2\times 3\times 4)\left(\frac{1}{2}+\frac{1}{3}+\frac{1}{4}\right) =$
$\text{(A)}\ 1 \qquad \text{(B)}\ 3 \qquad \text{(C)}\ 9 \qquad \text{(D)}\ 24 \qquad \text{(E)}\ 26$


	\item (\textit{AMC 8 1989}) There are $2$ boys for every $3$ girls in Ms. Johnson's math class.  If there are $30$ students in her class, what percent of them are boys?


 $\text{(A)}\ 12\% \qquad \text{(B)}\ 20\% \qquad \text{(C)}\ 40\% \qquad \text{(D)}\ 60\% \qquad \text{(E)}\ 66\frac{2}{3}\%$


	\item (\textit{AMC 8 1989}) What is the number of degrees in the smaller angle between the hour hand and the minute hand on a clock that reads seven o'clock?


 $\text{(A)}\ 50^\circ \qquad \text{(B)}\ 120^\circ \qquad \text{(C)}\ 135^\circ \qquad \text{(D)}\ 150^\circ \qquad \text{(E)}\ 165^\circ$


	\item (\textit{AMC 8 1989}) Which of the five "T-like shapes" would be symmetric to the one shown with respect to the dashed line?


 
                    \begin{figure}[H]
                        \centering
                        \includegraphics[width=0.5\textwidth]{Problem/AMC_8_AJHSME_1989_11_Problem_Figure_1.png}
                    \end{figure}
                    


	\item (\textit{AMC 8 1989}) $\frac{1-\frac{1}{3}}{1-\frac{1}{2}} =$
$\text{(A)}\ \frac{1}{3} \qquad \text{(B)}\ \frac{2}{3} \qquad \text{(C)}\ \frac{3}{4} \qquad \text{(D)}\ \frac{3}{2} \qquad \text{(E)}\ \frac{4}{3}$


	\item (\textit{AMC 8 1989}) $\frac{9}{7\times 53} =$
$\text{(A)}\ \frac{.9}{.7\times 53} \qquad \text{(B)}\ \frac{.9}{.7\times .53} \qquad \text{(C)}\ \frac{.9}{.7\times 5.3} \qquad \text{(D)}\ \frac{.9}{7\times .53} \qquad \text{(E)}\ \frac{.09}{.07\times .53}$


	\item (\textit{AMC 8 1989}) When placing each of the digits $2,4,5,6,9$ in exactly one of the boxes of this subtraction problem, what is the smallest difference that is possible?


 $\text{(A)}\ 58 \qquad \text{(B)}\ 123 \qquad \text{(C)}\ 149 \qquad \text{(D)}\ 171 \qquad \text{(E)}\ 176$

 \[\begin{tabular}[t]{cccc}  & \boxed{} & \boxed{} & \boxed{} \\ - & & \boxed{} & \boxed{} \\ \hline \end{tabular}\]


	\item (\textit{AMC 8 1989}) The area of the shaded region $BEDC$ in parallelogram $ABCD$ is


 
                    \begin{figure}[H]
                        \centering
                        \includegraphics[width=0.5\textwidth]{Problem/AMC_8_AJHSME_1989_15_Problem_Figure_1.png}
                    \end{figure}
                    
$\text{(A)}\ 24 \qquad \text{(B)}\ 48 \qquad \text{(C)}\ 60 \qquad \text{(D)}\ 64 \qquad \text{(E)}\ 80$


	\item (\textit{AMC 8 1989}) In how many ways can $47$ be written as the sum of two primes?


 $\text{(A)}\ 0 \qquad \text{(B)}\ 1 \qquad \text{(C)}\ 2 \qquad \text{(D)}\ 3 \qquad \text{(E)}\ \text{more than 3}$


	\item (\textit{AMC 8 1989}) The number $\text{N}$ is between $9$ and $17$.  The average of $6$, $10$, and $\text{N}$ could be


 $\text{(A)}\ 8 \qquad \text{(B)}\ 10 \qquad \text{(C)}\ 12 \qquad \text{(D)}\ 14 \qquad \text{(E)}\ 16$


	\item (\textit{AMC 8 1989}) Many calculators have a reciprocal key $\boxed{\frac{1}{x}}$ that replaces the current number displayed with its reciprocal.  For example, if the display is $\boxed{00004}$ and the $\boxed{\frac{1}{x}}$ key is pressed, then the display becomes $\boxed{000.25}$.  If $\boxed{00032}$ is currently displayed, what is the fewest positive number of times you must depress the $\boxed{\frac{1}{x}}$ key so the display again reads $\boxed{00032}$?


 $\text{(A)}\ 1 \qquad \text{(B)}\ 2 \qquad \text{(C)}\ 3 \qquad \text{(D)}\ 4 \qquad \text{(E)}\ 5$


	\item (\textit{AMC 8 1989}) The graph below shows the total accumulated dollars (in millions) spent by the Surf City government during $1988$.  For example, about $.5$ million had been spent by the beginning of February and approximately $2$ million by the end of April.  Approximately how many millions of dollars were spent during the summer months of June, July, and August?


 $\text{(A)}\ 1.5 \qquad \text{(B)}\ 2.5 \qquad \text{(C)}\ 3.5 \qquad \text{(D)}\ 4.5 \qquad \text{(E)}\ 5.5$

                    \begin{figure}[H]
                        \centering
                        \includegraphics[width=0.5\textwidth]{Problem/AMC_8_AJHSME_1989_19_Problem_Figure_1.png}
                    \end{figure}
                    


	\item (\textit{AMC 8 1989}) The figure may be folded along the lines shown to form a number cube.  Three number faces come together at each corner of the cube.  What is the largest sum of three numbers whose faces come together at a corner?


 
                    \begin{figure}[H]
                        \centering
                        \includegraphics[width=0.5\textwidth]{Problem/AMC_8_AJHSME_1989_20_Problem_Figure_1.png}
                    \end{figure}
                    
$\text{(A)}\ 11 \qquad \text{(B)}\ 12 \qquad \text{(C)}\ 13 \qquad \text{(D)}\ 14 \qquad \text{(E)}\ 15$


	\item (\textit{AMC 8 1989}) Jack had a bag of $128$ apples.  He sold $25\%$ of them to Jill.  Next he sold $25\%$ of those remaining to June.  Of those apples still in his bag, he gave the shiniest one to his teacher.  How many apples did Jack have then?


 $\text{(A)}\ 7 \qquad \text{(B)}\ 63 \qquad \text{(C)}\ 65 \qquad \text{(D)}\ 71 \qquad \text{(E)}\ 111$


	\item (\textit{AMC 8 1989}) The letters $\text{A}$, $\text{J}$, $\text{H}$, $\text{S}$, $\text{M}$, $\text{E}$ and the digits $1$, $9$, $8$, $9$ are "cycled" separately as follows and put together in a numbered list:

 \[\begin{tabular}[t]{lccc}  & & AJHSME & 1989 \\  & & & \\ 1. & & JHSMEA & 9891 \\ 2. & & HSMEAJ & 8919 \\ 3. & & SMEAJH & 9198 \\  & & ........ &  \end{tabular}\]


 What is the number of the line on which $\text{AJHSME 1989}$ will appear for the first time?


 $\text{(A)}\ 6 \qquad \text{(B)}\ 10 \qquad \text{(C)}\ 12 \qquad \text{(D)}\ 18 \qquad \text{(E)}\ 24$


	\item (\textit{AMC 8 1989}) An artist has $14$ cubes, each with an edge of $1$ meter.  She stands them on the ground to form a sculpture as shown.  She then paints the exposed surface of the sculpture.  How many square meters does she paint?


 $\text{(A)}\ 21 \qquad \text{(B)}\ 24 \qquad \text{(C)}\ 33 \qquad \text{(D)}\ 37 \qquad \text{(E)}\ 42$

                    \begin{figure}[H]
                        \centering
                        \includegraphics[width=0.5\textwidth]{Problem/AMC_8_AJHSME_1989_23_Problem_Figure_1.png}
                    \end{figure}
                    


	\item (\textit{AMC 8 1989}) Suppose a square piece of paper is folded in half vertically.  The folded paper is then cut in half along the dashed line.  Three rectangles are formed-a large one and two small ones.  What is the ratio of the perimeter of one of the small rectangles to the perimeter of the large rectangle?


 $\text{(A)}\ \frac{1}{2} \qquad \text{(B)}\ \frac{2}{3} \qquad \text{(C)}\ \frac{3}{4} \qquad \text{(D)}\ \frac{4}{5} \qquad \text{(E)}\ \frac{5}{6}$

                    \begin{figure}[H]
                        \centering
                        \includegraphics[width=0.5\textwidth]{Problem/AMC_8_AJHSME_1989_24_Problem_Figure_1.png}
                    \end{figure}
                    


	\item (\textit{AMC 8 1989}) Every time these two wheels are spun, two numbers are selected by the pointers.  What is the probability that the sum of the two selected numbers is even?


 $\text{(A)}\ \frac{1}{6} \qquad \text{(B)}\ \frac{3}{7} \qquad \text{(C)}\ \frac{1}{2} \qquad \text{(D)}\ \frac{2}{3} \qquad \text{(E)}\ \frac{5}{7}$

                    \begin{figure}[H]
                        \centering
                        \includegraphics[width=0.5\textwidth]{Problem/AMC_8_AJHSME_1989_25_Problem_Figure_1.png}
                    \end{figure}
                    


\end{enumerate}